\section{Introduction}

Assembling a functional cortex requires neurons to form specific synaptic connections and then refine them through activity-dependent mechanisms. Early circuit specification relies on surface recognition molecules that identify candidate partners, while subsequent refinement is driven by patterned spontaneous activity before sensory experience becomes available \citep{sanes2020synaptic}. Building on this developmental plan, a central organizational motif in the adult neocortex is the clustering of functionally similar synaptic inputs along short stretches of dendrite. In primary visual cortex of mice, the spatial receptive fields of synaptic inputs at neighboring spines are more similar than expected by chance, and co-tuned inputs are more likely to be co-located along individual branches \citep{iacaruso2017synaptic}. In ferret V1, clustering is organized around orientation preference and tracks the input-output nonlinearity of the parent cell \citep{wilson2016orientation}. Spontaneous network activity even in reduced preparations already reveals local synchronization between adjacent spines \citep{takahashi2012locally}.

Clustering matters because cortical pyramidal dendrites are electrically excitable and can sum coincident nearby inputs supralinearly. Radial oblique dendrites of CA1 pyramidal neurons generate NMDA-dependent dendritic spikes when synchronous inputs converge within tight temporal and spatial windows \citep{losonczy2006integrative}, and similar regenerative events in V1 and somatosensory cortex amplify stimulus responses and sharpen tuning \citep{smith2013dendritic,lavzin2012nonlinear}. Dendritic spines themselves add a substantial impedance that amplifies unitary synaptic inputs and encourages cooperativity between neighboring inputs \citep{harnett2012synaptic}. The pyramidal neuron behaves, in this view, as a two-layer network in which each branch computes a local nonlinear subunit \citep{poirazi2003pyramidal}, and the cortex exploits this design at a circuit scale \citep{larkum2013cellular}. Experimentally silencing dendritic nonlinearities in V1 or barrel cortex degrades tuning selectivity \citep{smith2013dendritic,lavzin2012nonlinear,wilson2018differential}, directly coupling subcellular synaptic geometry to perceptual computation.

Despite this strong case for clustered inputs in adult cortex, how clustering emerges during development remains unresolved. Three basic questions have not been answered at the resolution of individual functional synapses in vivo: when in development do inputs become clustered, whether synapses are deposited already clustered or are sorted into clusters after formation, and whether inputs distribute smoothly along dendrites or concentrate in spatially discrete domains. Earlier two-photon surveys of dendritic inputs to adult neurons \citep{jia2010dendritic} and reviews of developmental synaptic organization \citep{winnubst2012synaptic} identified clustering as a candidate organizing principle but could not address these questions directly. The critical developmental window is the second postnatal week in rodents, when synapse density rises rapidly and approaches mature levels just before eye opening at postnatal day 14.

To address this gap, we combined in vivo whole-cell voltage-clamp recordings with two-photon dendritic calcium imaging using GCaMP6s \citep{chen2013ultrasensitive} in layer 2/3 pyramidal neurons of mouse primary visual cortex between postnatal day 8 and 13. We used QX-314 in the pipette to block action potentials \citep{takahashi2012locally} and NMDA-facilitating holding potentials to resolve transmission at both spine and shaft synapses \citep{jia2010dendritic,svoboda1999spread}. Resonant scanning combined with piezo-driven focusing allowed us to sample much larger portions of the dendritic tree than in prior work, yielding 354 functional synapses and 3440 individual transmission events across 12 dendritic areas from 11 mice. Our measurements reveal that clustered inputs assemble into spatially and functionally distinct dendritic domains already before visual experience, and that local co-activity is correlated with synapse stabilization -- providing a developmental account of how the adult clustered input architecture arises.

\section{Related Work}

\paragraph{Clustered synaptic inputs in adult cortex.}
Imaging of dendritic spines in V1 has established that synapses carrying related visual features are spatially clustered along dendrites. \citet{iacaruso2017synaptic} mapped spatial receptive fields onto individual spines in mouse V1 and found that inputs from nearby visual locations preferentially cluster on neighboring spines, and \citet{wilson2016orientation} showed an analogous organization by orientation preference in ferret V1, with clustering magnitude tied to each cell's subthreshold nonlinearity. \citet{jia2010dendritic} first demonstrated that orientation-tuned synaptic calcium transients could be resolved in vivo in dendritic segments of mouse V1, and \citet{takahashi2012locally} revealed local synchronization of neighboring spines during spontaneous network activity. \citet{wilson2018differential} extended this framework to show that differential tuning of excitation and inhibition, rather than excitatory clustering alone, shapes direction selectivity.

\paragraph{Dendritic computation and its functional consequences.}
The functional relevance of clustered inputs depends on active dendritic integration. \citet{losonczy2006integrative} showed that hippocampal CA1 radial oblique dendrites sum synchronous inputs supralinearly via NMDA receptor spikes. \citet{harnett2012synaptic} quantified how spine neck impedance amplifies unitary inputs and promotes cooperativity between co-active inputs. In vivo, \citet{smith2013dendritic} linked dendritic spikes to enhanced orientation tuning in mouse V1 pyramidal cells, and \citet{lavzin2012nonlinear} found analogous NMDA-dependent amplification of angular tuning in barrel cortex layer 4. These mechanisms motivate the two-layer view of pyramidal function \citep{poirazi2003pyramidal} and the broader organizational hypothesis that active dendrites are a general cortical computing primitive \citep{larkum2013cellular,makara2013variable}.

\paragraph{Developmental assembly of cortical circuits.}
The formation of synapses between cortical neurons is guided by recognition molecules \citep{sanes2020synaptic} and is refined by patterned spontaneous activity during late embryonic and early postnatal life. In the developing mouse visual system, retinal waves propagate through the thalamus into primary visual cortex and are the dominant source of activity before eye opening \citep{ackman2012retinal}. \citet{siegel2012peripheral} identified two distinct patterns of V1 spontaneous activity in vivo whose participation rates are homogeneous across neurons, providing the input statistics that shape the first draft of cortical circuitry. Monosynaptic and quantal measurements of cortical connections \citep{markram1997physiology,silver2003high} further indicate that release-probability heterogeneity, rather than firing-rate diversity, determines the efficacy of individual inputs. Reviews have long proposed that clustering may arise through local activity-dependent plasticity \citep{winnubst2012synaptic}, and structural in vivo imaging demonstrates the remarkable dynamism of both excitatory spines and inhibitory contacts during circuit maturation \citep{villa2016inhibitory,kasai2023mechanical}. What has been missing, and what the present study provides, is a longitudinal map at single-synapse resolution of how functional inputs become clustered and organized into distinct dendritic domains during the second postnatal week.
